\begin{thema}{A}
\begin{erwthma}
    \item Να αποδείξετε ότι $ (x)'=1 $.
    \item Πότε η ευθεία $x=x_0$ ονομάζεται κατακόρυφη ασύμπτωτη της γραφικής παράστασης μιας συνάρτησης $f$;
    \item Πώς ορίζεται η αντίστροφη συνάρτηση $f^{-1}$ μιας $1-1$ συνάρτησης $f$;
    \item Να χαρακτηρίσετε τις προτάσεις που ακολουθούν, γράφοντας στο τετράδιό σας, δίπλα στο γράμμα που αντιστοιχεί σε κάθε πρόταση, τη λέξη \textbf{Σωστό}, αν η πρόταση είναι σωστή, ή \textbf{Λάθος}, αν η πρόταση είναι λανθασμένη.
    \begin{alist}
        \item Αν $f(x) > 0$ κοντά στο $x_0$, τότε $\lim\limits_{x\to x_0} f(x) > 0$.
        \item H $e^x$ δεν έχει κατακόρυφη ασύμπτωτη.
        \item Αν $f$ συνεχής και $f(a)f(\beta)=0$, υπάρχει $x_0 \in [a, \beta]$ με $f(x_0)=0$.
        \item Η εφαπτομένη της γραφικής παράστασης μιας συνάρτησης απαγορεύεται να τέμνει την καμπύλη σε άλλα σημεία.
        \item Κάθε σημείο $x_0$ με $f'(x_0)=0$ λέγεται κρίσιμο σημείο.
    \end{alist}
    \item Τo ολοκλήρωμα $\dint_1^e \frac{1}{x} dx$ ισούται με:
    \begin{tasks}[label=\let\textdexiakeraia\relax\alph*.](4)
        \task $e-1$
        \task $1$
        \task $\frac{1}{e}$
        \task $e$
    \end{tasks}
\end{erwthma}
\end{thema}
